\chapter{Resultados preliminares}
\label{chap:resultados-preliminares}

\section{Resultados do Motor de Dados (EDGAR)}
\label{sec:preliminares_edgar}

\begin{itemize}
    \item \textbf{Segmentação e Análise de Layout de Documentos (DLA):} Apresentação de uma figura demonstrando a aplicação prática do modelo de visão computacional sobre uma página de um artigo ou livro didático, evidenciando os delimitadores espaciais (\textit{bounding boxes}) isolando blocos de texto, tabelas e elementos gráficos.
    \item \textbf{Extração e Estruturação de Metadados:} Exibição do fragmento de código contendo o \textit{payload} JSON gerado pelo pipeline \textit{offline}, comprovando a capacidade do EDGAR de indexar o conteúdo textual associado às suas respectivas coordenadas geométricas de pixels e numeração de páginas.
\end{itemize}

\section{Protótipo da Interface Cliente (Luminia Reader)}
\label{sec:preliminares_luminia}

\begin{itemize}
    \item \textbf{Ambiente Integrado de Estudos:} Planejamento de figura contendo o protótipo de alta fidelidade da tela principal do leitor de PDF com o painel lateral de tutoria por IA de dados simulados.
    \item \textbf{Mecanismo de Destaque Geométrico:} Planejamento de figura demonstrando o comportamento visual da interface ao interpretar as coordenadas do JSON, aplicando o sombreamento colorido sobre o componente correspondente no documento para guiar o estudante.
\end{itemize}